%\documentclass[runningheads]{llncs}
\documentclass[11pt]{article}
\usepackage{amsfonts,amssymb}
\usepackage{plain}
\setcounter{tocdepth}{3}
\usepackage{color}
\usepackage{mdframed}
\usepackage{graphicx}
\usepackage{listings}
\usepackage{multicol}
%\usepackage{algorithmic,algorithm}
\usepackage{textcomp,booktabs}
\usepackage{graphicx,booktabs,multirow}
\usepackage{booktabs}
\newsavebox{\tablebox}
\usepackage{times}
\usepackage{hyperref}
\usepackage{ulem}

\topmargin -0.5cm \oddsidemargin 0cm \evensidemargin 0cm \textheight
22cm \textwidth 16cm




\renewcommand{\baselinestretch}{1.0}
 
\newcommand{\tb}{\textcolor{blue}}
\newcommand{\tr}{\textcolor{red}}

%-------------------------------------------------------------------------
\begin{document}

\title{Title of Project Proposal}


\author{
Name 1\\
Name 2\\
Name 3\\
Name 4\\
Name 5\\
}




\maketitle

\abstract{This project innovatively presents a city to tourists. Each Australian city has attractions including landmarks, buildings, monuments, etc. Currently tourists must browse the Internet or brochures to learn their stories. This is inconvenient during a tour, especially for foreigners. This calls for a convenient and instant information access mechanism. Considering everyone has a smart phone with camera, a mobile application can be developed with our image recognition research. By simply taking a picture of an attraction, tourists can instantly access its story, presented visually or verbally in their mother tongue. Functions such as photo-sharing and local advertisements can be added.

\tr{Fuchun: Put the paragraph of project description here. It should cover
\begin{itemize}
\item the background, 
\item the current solution and its disadvantages (that you focus on)
\item  and what this project will contribute
\end{itemize}
You must edit the proposal in this template, not other templates. Pls contact Fuchun if you don't know how to complete the second assignment.
}

\newpage

{\bf Statement of Contribution Weights:}

\begin{center}
\begin{tabular}{cc}\hline
Names & Weights \\ 
Name 1 &  100$\%$ \\
Name 2 &  100$\%$ \\
Name 3 &  100$\%$ \\
Name 4 &  90$\%$ \\
Name 5 &  100$\%$ \\
Name 6&  100$\%$ \\ \hline
\end{tabular}
\end{center}

\noindent{\bf  Explanations}: Put your brief explanations why Name 4 got 90 percents only.    For example,  this student missed many meetings and discussions. The number will be for reference only.    If your contribution weights are identical, no need to explain.

\newpage

\section{\tb{Introduction}}

\tr{
\begin{itemize}
\item The project proposal must have at least 10 pages including the title page but excluding references. 
\item You can insert some pics or tables to further explain what you are going to introduce. 
\item You must study how to cite references.  Note that part of contents of this proposal are allowed to be directly used in your final report.
\item You are not allowed to change the format or font size of this Latex template.
\item You can add other sections or subsections if needed.
\end{itemize}
}

\medskip


There are many ways of introducing the background of your project. Here, I introduce one template (You don't need to exactly follow this template):

The introduction can be seen as the extension of the abstract.
\begin{itemize}
\item Introduce the background of problem you are going to mention.

\item Introduce the problem (why this is important)


\item Briefly introduce and summarize related solutions, and identify the gaps. Put the detailed introductions in the related work.


\item Introduce the aims in the subsection in details.

\item  Introduce related work  in the subsection in details.
\end{itemize}


\subsection{Aims}

\subsection{Related Work}


Briefly introduce each solution, advantage and disadvantage of each solution.

TIPS: You can introduce details about each solution if it is hard for you to write up to 10 pages (from introduction section to the last section without counting references).


\section{\tb{Methodology}}
The approach/methodology section describes the methods and materials that you intend to use within your project. It will detail how you will go about acquiring the information you need and explain how the project will be managed. If a particular theory or model is to be used, this should be stated. You could also include details such as timelines,locations, or the software to be used.





You can additionally mention others, but you must at least mention the following points in each subsection.

\subsection{Project Challenges}
 \textbf{Authenticated Key Exchange}
 One of the main challenges of this project is implementing an authenticated key exchange mechanism. This involves two key difficulties. First, during the initial exchange of public keys, the system must ensure that the received public key genuinely belongs to the intended user rather than to a man-in-the-middle attacker. Without a reliable authentication mechanism, an attacker could intercept the communication and replace the users’ public keys with their own.
 Second, the initial version of the system will use a manual key exchange process, in which Alice sends her public key directly to Bob. Although this approach is relatively straightforward to implement, the ultimate goal is to automate the key exchange process while maintaining the authenticity and security of the exchanged public keys. Designing a secure and user-friendly transition from manual to automatic key exchange will therefore be an important technical challenge.

 To address these challenges, the system will use two levels of cryptographic keys: long-term identity keys and conversation-specific ephemeral keys. During the initial authentication process, a human-readable verification code will be generated from the fingerprint of each user’s long-term public key. Alice and Bob will compare this code through an out-of-band channel other than WhatsApp. This allows them to verify that the long-term public keys genuinely belong to the intended users. Certificate-based authentication may also be considered as an alternative if a suitable trust model is available.

 After the users’ long-term identities have been verified, the system will support automatic session-key establishment for individual WhatsApp conversations. When a user opens a conversation, the extension will obtain the conversation ID from the WhatsApp platform. This ID will be used to identify the conversation and to create or retrieve a conversation-specific ephemeral X25519 key pair. The extension will then automatically exchange the ephemeral public keys between the two participants.

 After receiving the other participant’s ephemeral public key, each party will use X25519 to calculate the same shared secret. HKDF will derive a temporary message-encryption key from this shared secret. The WhatsApp conversation ID, together with the participants’ identities and relevant protocol information, will be included in the key-derivation context. As a result, the derived message key will be cryptographically bound to the intended WhatsApp conversation and will not be directly reusable in another conversation.

 \textbf{Group Messaging Encryption}
 Another significant challenge is implementing end-to-end encryption for group conversations. Unlike a one-to-one conversation, a group conversation involves multiple participants, which makes key distribution and management more complex. The system must securely establish a shared group key and ensure that it is available only to authenticated group members.

 Group membership may change over time. When a new participant joins, the system must securely distribute the necessary keying material without allowing the new member to decrypt messages sent before they joined. When a participant leaves or is removed, a new group key must be generated and securely distributed to the remaining members so that the former member cannot decrypt future messages. Therefore, the system must support secure and efficient key rotation whenever group membership changes.

 The group key will be used to derive the temporary message-encryption key for the corresponding WhatsApp group conversation. The WhatsApp conversation ID and relevant protocol information will be incorporated into the key-derivation context to bind the derived key to the intended group. Designing a mechanism that supports authenticated membership, secure group-key distribution, membership changes, and automatic key updates will therefore be an important technical challenge.

 Due to its additional complexity, group messaging encryption may be implemented after the one-to-one messaging and key-establishment mechanisms have been successfully developed and evaluated.



 \subsection{Plan and Time Line}

~~


~\\
\section{\tb{Outcomes}}

This section should include the expected impacts and benefits of the proposed project. It is important that these be specific and realistic.

Example: “Outcomes of this project will include:

\begin{itemize}
\item a prototype of an engrossing, original computer game of professional standard

\item An academic paper that solves XXX problems.

\item for team members, valuable insights, enhanced skills and an original work for portfolios.”
\end{itemize}


\section{\tb{Conclusion}}
 
 
\tr{You can use any format of the reference as long as consistent. Pls search what is  Latex reference bib
}


\url{https://www.overleaf.com/learn/latex/Bibliography_management_with_bibtex}
 
 \begin{thebibliography}{9}
\bibitem{texbook}
Donald E. Knuth (1986) \emph{The \TeX{} Book}, Addison-Wesley Professional.

\bibitem{lamport94}
Leslie Lamport (1994) \emph{\LaTeX: a document preparation system}, Addison
Wesley, Massachusetts, 2nd ed.
\end{thebibliography}

\end{document}
